% !PRINTED 53
\weberpage{53}
%<<CONT

\[
\alpha_1 = \frac{\eta_1 a_1}{2},\ \alpha_2 = \frac{\eta_2 a_2}{2} \cdots \alpha_\lambda = \frac{\eta_\lambda a_\lambda}{2},\ \alpha_{\lambda+1}=0 \ldots \alpha_\nu = 0,
\]

in which $\eta_1, \eta_2 \ldots \eta_\lambda$ can be equal to 0 or equal to 1, and one obtains all two-sided elements in the form

\[
A_1^{\frac{\eta_1 a_1}{2}} A_2^{\frac{\eta_2 a_2}{2}} \ldots A_\lambda^{\frac{\eta_\lambda a_\lambda}{2}},
\]

and their number, since all combinations of 0 and 1 are admissible for the exponents $\eta$, is $= 2^\lambda$.

If $\chi$ is an arbitrary character of $S$, then, if $A$ is a two-sided element, $\chi(A) = \pm 1$, since

\[
\chi(A)^2 = \chi(1) = 1
\]

holds.

Similarly, we call a character two-sided if it is a two-sided element in the group of characters. Since the group of characters is isomorphic with the group $S$, there are just as many two-sided characters as two-sided elements, namely $2^\lambda$. By \S.\,11, (17) they are represented by:

\[
\chi_1^{\frac{\eta_1 a_1}{2}} \chi_2^{\frac{\eta_2 a_2}{2}} \ldots \chi_\lambda^{\frac{\eta_\lambda a_\lambda}{2}}.
\]

\gesperrt{The two-sided characters have the value $\pm 1$ for every element.}

The two-sided elements form by themselves a group $T$ of degree $2^\lambda$. Likewise the two-sided characters form a group isomorphic with it, and one obtains the group $U$ reciprocal to $T$ by seeking out all elements $A$ for which all two-sided characters have the value $+1$. This group, which by \S.\,12, 7. is a divisor of $S$ of index $2^\lambda$, depends, by the last definition, neither on the choice of the basis nor on the roots of unity representing the characters, and is completely determined by the nature of the group $S$. If we denote it by $G$, then the system of subsidiary groups

\weq{3}{G,\ G_1,\ G_2 \ldots G_{2^\lambda - 1}}

is characterised by the fact that for all elements of one of these systems the two-sided characters yield one and the same system of values. The systems $G, G_1, G_2 \ldots$ are called the \gesperrt{genera} (Genera) contained in $S$. The group $G$ in particular is called the \gesperrt{principal genus}.
